% ============================================================
% PHỤ LỤC A - GLOSSARY
% ============================================================

Phụ lục này được thêm để người mới không phải đoán nghĩa các từ tiếng Anh xuất hiện trong Master Plan.

\begin{longtable}{p{3.4cm} p{11.2cm}}
\toprule
\textbf{Thuật ngữ} & \textbf{Giải thích trong ngữ cảnh LabGuard} \\
\midrule
\endfirsthead
\toprule
\textbf{Thuật ngữ} & \textbf{Giải thích trong ngữ cảnh LabGuard} \\
\midrule
\endhead
Agent & Chương trình chạy trên máy cần giám sát. Nó thu trạng thái máy, gửi heartbeat/telemetry và reconnect khi cần. Không phải AI Agent. \\
Telemetry & Dữ liệu trạng thái được gửi từ xa: CPU, RAM, NIC, uptime, service status. \\
Heartbeat (HB) & Gói nhỏ gửi định kỳ để chứng minh Agent vẫn còn sống ở tầng ứng dụng. \\
Liveness & Câu hỏi ``thành phần này còn đang hoạt động không?''; thường được quyết định từ heartbeat + timeout/deadline. \\
Central Server & Máy chủ giữ nghiệp vụ chính: device state, rule, Incident, DB và realtime cho Admin. \\
Incident & Phiếu sự cố có vòng đời từ phát hiện đến xử lý xong; không đồng nghĩa một dòng log. \\
ACK / Acknowledge & Xác nhận đã nhận sự cố hoặc đã nhận message, tùy ngữ cảnh protocol/Incident. \\
TCP & Giao thức hướng kết nối cung cấp byte stream tin cậy, có thứ tự; không giữ biên message của ứng dụng. \\
Byte stream & TCP đưa ra chuỗi byte liên tục; một lần đọc có thể trả một phần message hoặc nhiều message dính nhau. \\
Framing & Quy tắc để receiver biết message bắt đầu/kết thúc ở đâu trên byte stream. \\
Length-prefix & Cách framing: gửi độ dài trước, rồi gửi đúng số byte payload. LabGuard dự kiến dùng 4 byte length. \\
Partial read & Một lần ReadAsync nhận ít byte hơn số byte của message; code phải đọc tiếp cho đủ. \\
Big-endian & Quy ước đặt byte có trọng số lớn trước khi biểu diễn số nhiều byte trên wire. \\
JSON / UTF-8 & Định dạng dữ liệu text và encoding dự kiến cho payload của protocol. \\
Envelope & Lớp bọc chung của message gồm version, messageType, requestId, deviceId, sentAtUtc và payload. \\
RequestId / CommandId & Mã định danh yêu cầu/lệnh để trace, retry và phát hiện trùng. \\
Deduplication & Loại bỏ việc xử lý lại một dữ liệu/lệnh đã commit trước đó. \\
Idempotency & Tính chất cho phép retry mà kết quả cuối không bị nhân đôi ngoài ý muốn. \\
Queue & Hàng đợi công việc/message giữa producer và consumer. \\
Bounded queue & Queue có giới hạn; khi đầy phải chờ/drop/fail theo policy thay vì ăn RAM vô hạn. \\
Backpressure & Cơ chế làm phía gửi chậm lại khi phía nhận không xử lý kịp. \\
Backoff & Sau lỗi kết nối, chờ một khoảng tăng dần trước khi thử lại để tránh reconnect storm. \\
Reconnect storm & Nhiều client reconnect đồng loạt, tạo tải đột biến cho server. \\
Rule engine & Khối đánh giá telemetry để quyết định có bất thường hay không. \\
Streak & Số mẫu liên tiếp cùng vi phạm; dùng để tránh phản ứng với một spike đơn lẻ. \\
Hysteresis & Dùng ngưỡng mở và ngưỡng hồi phục khác nhau để tránh trạng thái bật/tắt liên tục gần threshold. \\
ConditionState & Trạng thái kỹ thuật của điều kiện lỗi (ACTIVE/RECOVERED/UNKNOWN), tách khỏi workflow xử lý của con người. \\
Race condition & Kết quả sai/không ổn định do nhiều thao tác concurrent cùng đọc/ghi state. \\
Transaction & Nhóm thao tác DB được coi như một đơn vị nguyên tử: hoặc commit đủ, hoặc rollback. \\
rowversion & Datatype riêng của SQL Server (8 byte, do CSDL tự sinh khi row được insert/update; binary(8) thông thường không có semantics tự cập nhật này; không phải date/time); dùng làm token OCC. \\
UI Marshal & Chuyển giao dữ liệu từ luồng nền/mạng về luồng giao diện qua \texttt{Control.Invoke}, \texttt{BeginInvoke} hoặc \texttt{InvokeAsync} để bảo đảm an toàn đa luồng trên WinForms. \\
Snapshot & Ảnh trạng thái đầy đủ tại một thời điểm. \\
Delta / Event & Chỉ phần thay đổi sau snapshot; dùng để update realtime tiết kiệm hơn. \\
Slow-client isolation & Không để một Admin nhận dữ liệu chậm làm block mọi subscriber khác. \\
TLS & Lớp bảo vệ kênh truyền: mã hóa + xác thực endpoint theo certificate/config. \\
Certificate & Chứng thư số dùng trong TLS để xác nhận identity của endpoint và thiết lập kênh bảo mật. \\
Hashing mật khẩu & Biến password thành giá trị một chiều để server không lưu password gốc. \\
Salt & Dữ liệu ngẫu nhiên thêm vào password trước KDF để chống precomputed hash/rainbow table. \\
KDF & Password Key Derivation Function; cố tình tốn tính toán hơn hash nhanh để tăng chi phí brute force. \\
Authentication & Xác định ``ai đang kết nối''. \\
Authorization & Xác định identity đó được phép thao tác gì. \\
RBAC & Role-Based Access Control; phân quyền dựa trên 3 vai trò: Admin (toàn quyền, phân công, quản trị), Operator (xử lý sự cố: ACK/START/RESOLVE), Viewer (chỉ đọc). \\
Retention & Quy định dữ liệu giữ bao lâu trước khi aggregate/purge. \\
Purge & Xóa dữ liệu cũ theo retention policy. \\
Gateway & Điểm vào đứng trước nhiều Collector; ở P1 chọn node và relay kết nối. \\
Collector & Server Agent-facing ở P1, nhận/parse/auth/queue rồi chuyển dữ liệu lên Central. \\
Load balancing (LB) & Phân phối kết nối mới giữa nhiều server để tránh một node ôm hết tải. \\
Least connections & Chọn server đang có ít connection active hơn. \\
Health check & Kiểm tra node còn hoạt động và phản hồi đúng không. \\
Readiness & Cho biết node có sẵn sàng nhận công việc/kết nối mới không. \\
Failover & Khi node lỗi, client/workload chuyển sang node khác. Blueprint chấp nhận có gián đoạn. \\
Fencing & Ngăn session/node cũ tiếp tục ghi dữ liệu như thể nó vẫn còn quyền sau reconnect/failover. \\
Session generation & Số thế hệ của phiên; phiên reconnect mới phải mới hơn để loại update từ phiên cũ. \\
DAG & Directed Acyclic Graph; đồ thị phụ thuộc task không có vòng lặp. \\
WIP & Work In Progress; giới hạn số việc một người đang làm cùng lúc để tránh mở quá nhiều task. \\
MVP & Phiên bản nhỏ nhất nhưng hoàn chỉnh end-to-end, có giá trị demo và test thật. \\
P0/P1/P2/P3 & Mức ưu tiên: bắt buộc / nâng cao / stretch / tương lai. \\
Definition of Done & Điều kiện để task thực sự DONE: code + test + docs + review + owner hiểu và vấn đáp được. \\
Regression test & Test chạy lại để chứng minh sửa một lỗi/tính năng không làm hỏng hành vi đã đúng trước đó. \\
Failure injection & Chủ động tạo lỗi như disconnect/service down để kiểm tra hệ thống xử lý đúng. \\
Simulator & Chương trình tạo nhiều Agent/socket giả lập về thiết bị nhưng kết nối mạng và protocol là thật. \\
CGNAT & Cơ chế NAT của nhà mạng khiến thiết bị không có public IPv4 riêng, có thể làm inbound Internet khó/không khả thi. \\
\bottomrule
\end{longtable}
